\scanpage{176}\mspage{121}
ce qui s'écrit \unsure{ici|très} élégamment en termes d'homologie relative, par la condition que $\forall a\in\unsure{\mathcal A|A}=\mathbb D(e)$, la flèche induite par $f$
\diagram{\begin{tikzcd}[ampersand replacement=\&,column sep=large,row sep=small]
H_{\unsure{\bullet|\cdot}}^{\mathbb D}(X/S,a_X) \arrow[r] \arrow[d,equal,"\mathrm{def}"'] \& H_{\bullet}^{\mathbb D}(\deleted{\unsure{X}}Y/S,a_Y) \arrow[d,equal,"\mathrm{def}"] \\
\deleted{\unsure{p}}{p_X}_{!}(a_X) \& (p_Y)_{!}(a_Y)
\end{tikzcd}}
soit iso. (Utiliser le ``calcul'' des faisceaux d'homologie relative $(p_X)_!(a_X)$, $(p_Y)_!(a_Y)$, par la caractérisation des fibres de Der 4 bis).

Or, les changements de base $\mathbb D$-\uline{propres} $S'\longrightarrow S$ commutent aux $p_!$, de sorte qu'on trouve le

\uline{Lemme} La notion de \deleted{\unsure{$W$}}-équivalence $\mathbb D$-colocale sur $S$ dans $\mathrm{Cat}_{/S}$ est stable par \inserted{\unsure{tout}} changement de base $S'\longrightarrow S$ $\mathbb D$-propre. (Quel que soit le dérivateur $\mathbb D$)

\deleted{Cela} On se rappellera que les $ {}_{s\backslash}S\longrightarrow S$ sont justement des changements de base propres ! Ceci \unsure{peut|veut} dire que l'équivalence $\mathbb D$-colocale peut se décrire en exigeant que $f:X\longrightarrow Y$ soit transformé en $W$-équivalence par certains types de chgt. de base $S'\longrightarrow S$, en prenant au choix l'un des trois types

a) les chgts de base $ {}_{s\backslash}S\longrightarrow S$

(ce qui redonne la définition que nous avions \unsure{prise|posée})

b) les chgts de base qui sont des co-isom. locaux

c) les chgt. de base propres

d) les chgts de base $\mathbb D$-propres.

On ne peut appliquer ces résultats au \unsure{dériv.-|dériva-}
% archive sheet number: 175
